% M14_audit.tex -- publication, data-ownership and data-protection audit (state: repository commit 3fe6627 + this
% monograph's commits; checked 2026-10-14). Sources: research_record/00_PROJECT_CHARTER.md (sections 8, 9),
% 01_CURRENT_STATE.md (B, H.6), PUBLIC_EXCLUSIONS.md, 04_TASK_QUEUE.yaml (HUM items), data/provenance_*.json,
% git ls-files of the public checkout. Nothing was deleted, pushed or published by this audit.
\chapter{Publication, data-ownership and data-protection audit}
\label{ch:mono-audit}

Before any result of this project is published, three questions must be answered for every item it uses: who owns it,
whether it may be reproduced and redistributed, and whether it was already published by someone else. This chapter
records the state of these questions as found in the repository on 14 October 2026. It reports findings and
recommendations; the decisions belong to the humans named in the charter (section~9), and nothing was changed, deleted,
pushed or published in the course of the audit.

% =====================================================================================================================
\section{People, roles and authority}
\label{sec:mono-audit-people}

\begin{itemize}
\item \emph{Research owner:} Mahmoud Elrasheedy (TCAD and DFT; owner of the record). Mentor: Mukul Kumar; principal
investigator: Prof.\ Chao-Hsin Wu.
\item \emph{Fabrication group and data owner:} Prof.\ Min-Hung Lee's group, National Taiwan University (Dr.\ Mochamad
Januar, Zhao-Feng Luo); co-author of the published paper: Prof.\ Kou-Chen Liu (Chang Gung University).
\item \emph{AI agents:} Claude Code (computational execution, evidence engineering, record keeping) and Codex
(independent review). Under the charter, AI agents may not put claims into a manuscript, release data or push to a new
remote without explicit human approval; agreement between AI systems is not evidence. The charter itself is marked
\emph{proposed}; its adoption (\dlog{D-10}) is pending.
\end{itemize}

\section{Inventory by origin and status}
\label{sec:mono-audit-inventory}

\begin{table}[htbp]\centering\footnotesize
\caption{Items used by the project, their origin and their status. ``Public'': tracked in the public repository at the
audited commit.}
\label{tab:mono-audit-inventory}
\begin{tabularx}{\linewidth}{@{}>{\raggedright\arraybackslash}p{3.6cm}>{\raggedright\arraybackslash}p{2.5cm}>{\raggedright\arraybackslash}p{1.2cm}>{\raggedright\arraybackslash}X@{}}
\toprule
Item & Origin, owner & Public & Status and constraint \\\midrule
Four measured transfer curves (workbook and CSV; a second workbook copy under \file{inputs/astra_.../audit/sources/}) &
fabrication group & yes & devices of a published series (2--13~nm; the 31.8~nm curve was supplied with the data); the paper makes data available on reasonable request \citep[p.~10]{Januar2026}; permission for
public hosting, reuse in publications and AI processing \textbf{not recorded} (charter section~8; task HUM-001) \\
Published values and figures of \citet{Januar2026} & fabrication group; CC~BY article \citep[p.~1]{Januar2026} & yes (article) & may be cited
and, with attribution, reproduced; digitised values are our derived data \\
Supporting information of the paper & fabrication group & no & not in the project files \\
Unpublished fab data (stress, temperature, annealing, \InO{} devices, C--V) & fabrication group & no & not
available (\dlog{D-11}); would need consent before entering any AI tool \\
Crystalline DFT inputs, outputs and analyses & author (CERN and local compute) & yes & author-generated; releasable \\
Amorphous DFT inputs, specs, scripts & author & yes & releasable \\
Amorphous DFT outputs & author & no & blind; release blocked (\dlog{D-43}); excluded from the public repository \\
ATLAS decks and outputs, runs 0001--0074 & author & yes & author-generated; releasable \\
ATLAS runs 0075--0077, campaign C10\_DIAGb & author & no & sealed (state-only rule) \\
Full texts of 11 cited papers (\file{report_latex_full/digest/papers/}) & publishers and authors & \textbf{yes} &
copyright; several are not open access; tracked since 5 October 2026 although the \file{.gitignore} excludes the
equivalent \file{inputs/papers_text/} \\
ATLAS manual text (\file{report_latex_full/tools/atlas_manual_fulltext.json}, \file{atlas_manual_pages.txt}, about
3.5~MB each; page extracts in \file{digest/atlas_manual_pages.md}) & Silvaco (licensed software documentation) &
\textbf{yes} & copyright; \file{PUBLIC_EXCLUSIONS.md} excludes a 740-page manual text but not these files (contents not
opened by this audit) \\
Codex reviews, work chat, register & author and AI agents & partly & the work chat mixes reviews with sealed-adjacent
content; the register and decision log are public \\\bottomrule
\end{tabularx}
\end{table}

\section{Findings}
\label{sec:mono-audit-findings}

\begin{enumerate}
\item \textbf{Collaborator data in a public repository without recorded permission.} The four measured curves are the
fabrication group's data. They are in the public repository in two places. The paper makes data available on reasonable
request, which is not a licence to redistribute. The charter records this as a known issue awaiting the meeting with the
fabrication group (task HUM-001).
\item \textbf{Copyrighted full texts in the public repository.} The digest of eleven papers and the ATLAS manual text
files are publicly tracked. Removing them from the current tree would not remove them from the git history.
\item \textbf{Publication overlap.} The device data, an analytic model of them, and crystalline DFT of \InO{}
and W:\InO{} (PBE) with interface slabs are already published \citep{Januar2026}. Any paper of
this project that shows the curves or re-analyses the paper's metrics overlaps with that publication and needs the
fabrication group's agreement; a paper built on the author-generated calculations (the confinement series and its
analysis, the vacancy test, the amorphous protocol) overlaps less but still relies on the devices for motivation and
comparison.
\item \textbf{Original against re-analysed results.} Original to this project: the crystalline confinement series and
its laws, masses and alignment; the W-cell analysis and its correction; the vacancy test H-L; the amorphous protocol and
its numerical qualification; the TCAD models and campaign; the registered extraction and its reconciliation with the
published metrics. Re-analysed: every device metric (derived from the group's curves).
\item \textbf{AI-tool processing.} The charter requires the fabrication group's consent before unpublished fab data are
given to AI services. The four curves have been processed by AI agents throughout the project, including this
monograph; three of them belong to published devices, and all four were obtained as data, and the consent question for them is part of the
pending meeting.
\item \textbf{Timestamps of registrations.} Registrations are timestamped by commits and pushed tags in the repositories;
the charter's preferred external timestamp (a DOI from Zenodo or OSF) has not been obtained (paper-draft item P27).
\item \textbf{Sealed and blind material.} The audit found the blind guard of the \file{.gitignore} and
\file{PUBLIC_EXCLUSIONS.md} consistent with the absence of sealed runs and amorphous outputs in this checkout. One public
commit message contains sealed-adjacent information; the monograph does not use it.
\end{enumerate}

\section{Recommendations for the humans}
\label{sec:mono-audit-reco}

These are recommendations, not decisions; each needs the approval named in the charter.
\begin{enumerate}
\item Hold the meeting with the fabrication group (HUM-001) with a written agenda: permission for the curves already
public (keep, restrict or remove), permission and attribution for their use in publications, consent for AI processing,
data custody for future holdout data, and authorship and contribution roles for any paper.
\item Decide whether the copyrighted full texts and manual extracts should leave the public repository; if so, decide
whether the history should be rewritten (a destructive operation on a public repository that affects every clone).
\item Obtain an external timestamp (DOI) for the frozen register and the read-out freezes before any registered result is
reported outside the group.
\item Adopt or amend the charter (\dlog{D-10}), because the publication rules of this audit depend on it.
\item Keep the monograph and the manuscript scaffold internal until items 1 and 3 are settled.
\end{enumerate}

\section{Data and code availability, as they could be stated today}
\label{sec:mono-audit-availability}

\emph{Data.} The measured transfer curves belong to the fabrication group and are available from them on reasonable
request \citep[p.~10]{Januar2026}. All first-principles inputs and outputs and all device-simulation decks and outputs generated
in this work are available in the project repository, except calculations whose results are under a registered
blind-analysis protocol, which will be released when the protocol's release conditions are met. \emph{Code.} Structure
generators, workflow scripts, the extraction code, the frozen scorer and read-out package (with their freeze tags), and
the analysis scripts of this monograph are in the repository. Both statements are drafts for the humans to approve.

\begin{keyresult}
The project's own calculations are releasable, except the sealed and blind ones. The measured curves are a
collaborator's data, already public without recorded permission; copyrighted full texts of papers and of the ATLAS
manual are also public. Publication of any result that uses the curves needs the fabrication group's agreement on
permission, attribution, AI processing and authorship, which the pending meeting is meant to settle.
\end{keyresult}
